\documentclass[11pt]{article}
\usepackage[margin=1.2in]{geometry}
\pagestyle{empty}
\setlength{\parskip}{0.7\baselineskip}
\setlength{\parindent}{0pt}
\begin{document}
Dear Editors,

Please consider ``An Empirical Study of the Conformal Information Gap'' for publication in the International Journal of Forecasting.

The paper gives forecasting practice an exact price for a common shortcut: pooling residuals into a single empirical law before conditional structure has been removed. Under logarithmic loss that price is a conditional mutual information, and the paper's contribution is to measure it. A nested ladder of online transformations is scored prequentially on 701 economic series under proper scoring rules, in the sharpness-subject-to-calibration tradition. Conditional scale emerges as the largest and most reliable component of the measured gain, and the finding survives a matched-predictor comparison, a simulation of the coverage certificate under its own assumptions, and replication on an out-of-development holdout. The practical guidance for forecasters is direct: transform until residuals are conditionally homogeneous, pool in those coordinates, and add a conformal certificate afterward when one is needed.

A short companion note, ``Marginally Useful: Formalizing the Information Gap in Conformal Prediction,'' which proves the single-pool special case of the identity used here, is currently under review at The American Statistician. The present paper does not depend on that note's acceptance: it generalizes the identity to arbitrary retained representations, proves the transform-pooling theorem, and contributes the 701-series measurement study, none of which appear in the companion.

A preprint of this manuscript is posted on SSRN, as permitted by journal policy. The manuscript is not under consideration elsewhere. The anonymized manuscript withholds author identity, the companion self-citation, and code repository links; a public software library is cited in the third person in accordance with double-blind convention.

\vspace{6pt}
Thank you for your consideration.\\
Peter Cotton peter.cotton@microprediction.com
\end{document}
